% s-315-p315-m002; original PDF 315
\Aphorism{s-315-p315-m002}{369}
Was hat die abendländische Religiosität, nicht zu reden von der abendländischen Freigeisterei, an Tiefe und Kühnheit jener indischen Auffassung an die Seite zu stellen, dass die Götter die Andachtsvertiefung des Menschen \Sperr{fürchten}, der sich dadurch über die Natur und über sie selbst erhebt, und darum sein Busswerk durch Versuchungen stören! Man stelle dieses Ungethüm von Gedanken einmal in unser Europa hinein; wie zwergenhaft und halbwüchsig nimmt sich dagegen unsere Theolatrie aus, mit der man in dieser und jener Welt ein \Foreign{non sine dis animosus infans} bleibt! Wenn irgend etwas, so beweist solch ein Stück asiatischer Frömmigkeit, wie sehr unter den Befriedigungen des religiösen Bedürfnisses gerade die monotheistische in die Kinderstube gehört. Ohne ein wenig Hybris, die mit den Göttern auf gleichem Fusse verkehrt, wird Religion doch gar zu salzlos; das geprügelte Kind, das die Ruthe küsst, ist als Mann, als gläubiger Mensch keine erfreuliche Erscheinung.
\par

